\subsection{\texorpdfstring{Relation between exact sequences and elements of Ext \(_{A}\) (M, N)}{Relation between exact sequences and elements of Ext \_\{A\} (M, N)}}

THEOREM 1. --- Let n be an integer ≥1, and M and N two A-modules.

\begin{enumerate}
\def\labelenumi{\alph{enumi})}
\item
  Every element of Ext \(_{A}^{n}\) (M, N) is the class of an exact sequence (X, p.~117, def. 1).
\item
  Let \(0 \to \mathbf{N} \xrightarrow{f_{n+1}} \mathbf{R}_n \xrightarrow{f_n} \ldots \to \mathbf{R}_1 \xrightarrow{f_1} \mathbf{M} \to 0\)
\end{enumerate}

and

\[
0 \rightarrow \mathrm{N} \xrightarrow {f _ {n + 1} ^ {\prime}} \mathrm{R} _ {n} ^ {\prime} \xrightarrow {f _ {n} ^ {\prime}} \dots \rightarrow \mathrm{R} _ {1} ^ {\prime} \xrightarrow {f _ {1} ^ {\prime}} \mathrm{M} \rightarrow 0
\]

of exact sequences, \(\theta\) and \(\theta'\) the associated classes. The following conditions are equivalent:

\begin{enumerate}
\def\labelenumi{(\roman{enumi})}
\item
  \(\theta = \theta^{\prime}\) ;
\item
  there exists a commutative diagram with exact rows:
\end{enumerate}

\[
\begin{array}{c} 0 \longrightarrow \mathrm{N} \xrightarrow {f _ {n + 1}} \mathrm{R} _ {n} \longrightarrow \dots \longrightarrow \mathrm{R} _ {1} \xrightarrow {f _ {1}} \mathrm{M} \longrightarrow 0 \\ 0 \longrightarrow \mathrm{N} \longrightarrow \mathrm{R} _ {n} ^ {\prime \prime} \longrightarrow \dots \longrightarrow \mathrm{R} _ {1} ^ {\prime} \longrightarrow \mathrm{M} \longrightarrow 0 \\ 0 \longrightarrow \mathrm{N} \xrightarrow {f _ {n + 1} ^ {\prime}} \mathrm{R} _ {n} ^ {\prime} \longrightarrow \dots \longrightarrow \mathrm{R} _ {1} ^ {\prime} \xrightarrow {f _ {1} ^ {\prime}} \mathrm{M} \longrightarrow 0; \end{array}
\]

\begin{enumerate}
\def\labelenumi{(\roman{enumi})}
\setcounter{enumi}{2}
\tightlist
\item
  there exists a commutative diagram with exact rows:
\end{enumerate}

\[
\begin{array}{c} 0 \longrightarrow \mathrm{N} \xrightarrow {f _ {n + 1}} \mathrm{R} _ {n} \longrightarrow \dots \longrightarrow \mathrm{R} _ {1} \xrightarrow {f _ {1}} \mathrm{M} \longrightarrow 0 \\ 0 \longrightarrow \mathrm{N} \longrightarrow \mathrm{R} _ {n} ^ {\prime \prime} \longrightarrow \dots \longrightarrow \mathrm{R} _ {1} ^ {\prime \prime} \longrightarrow \mathrm{M} \longrightarrow 0 \\ 0 \longrightarrow \mathrm{N} \xrightarrow {f _ {n + 1} ^ {\prime}} \mathrm{R} _ {n} ^ {\prime} \longrightarrow \dots \longrightarrow \mathrm{R} _ {1} ^ {\prime} \xrightarrow {f _ {1} ^ {\prime}} \mathrm{M} \longrightarrow 0. \end{array}
\]

Let us prove \(a\) ). Let \(\alpha \in \operatorname{Ext}_{\mathbf{A}}^{n}(\mathbf{M}, \mathbf{N})\) , and let \(\mathbf{P}\) be a projective resolution of \(\mathbf{M}\) . Let \(a: \mathrm{P}(n) \to \mathbf{N}\) be a morphism of complexes representing \(\alpha\) ; its unique nonzero component is an A-homomorphism \(u: \mathrm{P}_{n} \to \mathbf{N}\) which satisfies \(u \circ d_{n+1} = 0\) , hence factors through \(u = \overline{u} \circ \delta_{n}\) , where \(\delta_{n}: \mathrm{P}_{n} \to \dot{\mathbf{Z}}_{n-1}\) is the map induced by \(d_{n}\) (we set \(Z_{n-1} = \operatorname{Im} d_{n}\) ) and \(\overline{u}\) is an A-homomorphism of \(Z_{n-1}\) to \(\mathbf{N}\) . By Remark 1, p.~117, the class \(\theta \in \operatorname{Ext}_{\mathbf{A}}^{n}(\mathbf{M}, Z_{n-1})\) of the exact sequence

\[
0 \rightarrow Z _ {n - 1} \rightarrow P _ {n - 1} \rightarrow \dots \rightarrow P _ {0} \rightarrow M \rightarrow 0
\]

is equal to the homotopy class of the morphism \((-1)^{n(n + 1) / 2}\delta_n\) . We therefore have

\[
\alpha = (- 1) ^ {n (n + 1) / 2} \overline {{{u}}} \circ \theta ,
\]

which, according to Cor. 3, p.~120, allows one to represent α as the class of an exact sequence.

It follows from cor. 1 of X, p.~120 that (ii) \(\Rightarrow\) (i) and that (iii) \(\Rightarrow\) (i). Suppose (i) is satisfied, and let P be a projective resolution of M. There exists a commutative diagram

\[
\begin{array}{c} 0 \longrightarrow \mathrm{N} \xrightarrow {f _ {n + 1}} \mathrm{R} _ {n} \longrightarrow \mathrm{R} _ {n - 1} \longrightarrow \dots \longrightarrow \mathrm{R} _ {1} \longrightarrow \mathrm{M} \longrightarrow 0 \\ u _ {n} \uparrow \quad \uparrow u _ {n - 1} \quad \uparrow u _ {n - 2} \quad \uparrow \quad \uparrow_ {\mathrm{I} _ {\mathrm{M}}} \\ \mathrm{P} _ {n} \xrightarrow {d _ {n}} \mathrm{P} _ {n - 1} \longrightarrow \mathrm{P} _ {n - 2} \longrightarrow \dots \longrightarrow \mathrm{P} _ {0} \longrightarrow \mathrm{M} \longrightarrow 0 \\ u _ {n} ^ {\prime} \downarrow \quad \downarrow u _ {n - 1} ^ {\prime} \quad \downarrow u _ {n - 2} ^ {\prime} \quad \downarrow \quad \downarrow_ {\mathrm{I} _ {\mathrm{M}}} \\ 0 \longrightarrow \mathrm{N} \xrightarrow {f _ {n + 1}} \mathrm{R} _ {n} ^ {\prime} \longrightarrow \mathrm{R} _ {n - 1} ^ {\prime} \longrightarrow \dots \longrightarrow \mathrm{R} _ {1} ^ {\prime} \longrightarrow \mathrm{M} \longrightarrow 0. \end{array}
\]

The morphisms of \(\mathrm{P}(n)\) to \(\mathbf{N}\) defined by \(u_{n}\) and \(u_{n}^{\prime}\) are homotopic, since they both belong to the class \((-1)^{n(n + 1) / 2}\theta\) , hence \(u_{n}^{\prime} - u_{n}\) is of the form \(w\circ d_n\) , where \(w:\mathrm{P}_{n - 1}\to \mathrm{N}\) is an A-homomorphism. By replacing \(u_{n - 1}^{\prime}\) by \(u_{n - 1}^{\prime} - f_{n + 1}^{\prime}\circ w\) and \(u_{n}^{\prime}\) by \(u_{n}\) one reduces to the case where \(u_{n} = u_{n}^{\prime}\) . This allows one to construct a new commutative diagram with exact rows:

\[
\begin{array}{c c c c c c c} 0 \longrightarrow \mathrm{N} \longrightarrow \mathrm{R} _ {n} \longrightarrow \mathrm{R} _ {n - 1} \longrightarrow \dots \longrightarrow \mathrm{R} _ {1} \longrightarrow \mathrm{M} \longrightarrow 0 \\ 1 _ {\mathrm{N}} \uparrow & \uparrow v & \uparrow u _ {n - 2} & \uparrow & \uparrow^ {1 _ {\mathrm{M}}} \\ 0 \longrightarrow \mathrm{N} \longrightarrow \mathrm{N} ^ {\prime} \longrightarrow \mathrm{P} _ {n - 2}. \longrightarrow \dots \longrightarrow \mathrm{P} _ {0} \longrightarrow \mathrm{M} \longrightarrow 0 \\ 1 _ {\mathrm{N}} \downarrow & \downarrow v ^ {\prime} & \downarrow u _ {n - 2} ^ {\prime} & \downarrow & \downarrow^ {1 _ {\mathrm{M}}} \\ 0 \longrightarrow \mathrm{N} \longrightarrow \mathrm{R} _ {n} ^ {\prime} \longrightarrow \mathrm{R} _ {n - 1} ^ {\prime} \longrightarrow \dots \longrightarrow \mathrm{R} _ {1} ^ {\prime} \longrightarrow \mathrm{M} \longrightarrow 0 \end{array}
\]

where \(\mathbf{N}'\) is the quotient of \(\mathbf{P}_{n-1} \oplus \mathbf{N}\) by the submodule formed by the pairs \((d_n(x), -u_n(x))\) for \(x \in \mathbf{P}_n\) , and where \(v\) (resp. \(v'\) ) is defined by passing to the quotient from the map \(u_{n-1} \oplus f_{n+1}\) (resp. \(u_{n-1}' \oplus f_{n+1}'\) ). Condition (ii) is therefore satisfied.

Suppose again that condition (i) is satisfied, and let E be an injective resolution of N. There exists a commutative diagram

\[
\begin{array}{c c c c c c c} 0 \longrightarrow \mathrm{N} \longrightarrow \mathrm{R} _ {n} \longrightarrow & \dots \longrightarrow \mathrm{R} _ {1} \longrightarrow \mathrm{M} \longrightarrow & 0 \\ 1 _ {\mathrm{N}} \Big \downarrow & \Big \downarrow v _ {0} & v _ {n - 1} \Big \downarrow & \Big \downarrow v _ {n} \\ 0 \longrightarrow \mathrm{N} ^ {\prime} \longrightarrow \mathrm{E} ^ {0} \stackrel {{\delta^ {0}}} {{\longrightarrow}} & \dots \longrightarrow \mathrm{E} ^ {n - 1} \stackrel {{\delta^ {n - 1}}} {{\longrightarrow}} \mathrm{E} ^ {n} \\ 1 _ {\mathrm{N}} \Big \uparrow & \Big \uparrow v _ {0} ^ {\prime} & v _ {n - 1} ^ {\prime} \Big \uparrow & \Big \uparrow v _ {n} ^ {\prime} \\ 0 \longrightarrow \mathrm{N} \longrightarrow \mathrm{R} _ {n} ^ {\prime} \longrightarrow & \dots \longrightarrow \mathrm{R} _ {1} ^ {\prime} \longrightarrow & \mathrm{M} \longrightarrow & 0 \end{array}
\]

and one shows as above that one may assume \(v_{n}^{\prime}=v_{n}\) . We then have a commutative diagram with exact rows

\[
\begin{array}{c c c c c c c c} 0 \longrightarrow N \longrightarrow R _ {n} \longrightarrow \dots \longrightarrow R _ {2} \longrightarrow R _ {1} \longrightarrow M \longrightarrow 0 \\ \downarrow^ {1 _ {N}} & \downarrow & & \downarrow & \downarrow & \downarrow^ {1 _ {M}} \\ 0 \longrightarrow N \longrightarrow E ^ {0} \longrightarrow \dots \longrightarrow E ^ {n - 2} \longrightarrow M ^ {\prime} \longrightarrow M \longrightarrow 0 \\ \uparrow^ {1 _ {N}} & \uparrow & & \uparrow & \uparrow & \uparrow^ {1 _ {M}} \\ 0 \longrightarrow N \longrightarrow R _ {n} ^ {\prime} \longrightarrow \dots \longrightarrow R _ {2} ^ {\prime} \longrightarrow R _ {1} ^ {\prime} \longrightarrow M \longrightarrow 0 \end{array}
\]

with \(\mathbf{M}' = \mathbf{M} \times_{\mathrm{R}_{1}} \mathbf{E}_{n-1}\) (cf. X, p. 120, cor. 3, b)). Condition (iii) is therefore satisfied, which completes the proof of the theorem.

Remark 1. --- If the ring A is Noetherian, and if the A-modules M and N are finitely generated, it follows from the proof of a) that every element of Ext \(_{A}^{n}\) (M, N) is the class associated with an exact sequence 0 → N → R \(_{n}\) → \ldots{} → R \(_{1}\) → M → 0 where the R \(_{i}\) are finitely generated.

COROLLARY. --- Let \(0 \to \mathbf{N} \xrightarrow{f} \mathbf{R} \xrightarrow{g} \mathbf{M} \to 0\) and \(0 \to \mathbf{N} \xrightarrow{f'} \mathbf{R}' \xrightarrow{g'} \mathbf{M} \to 0\) be two exact sequences, \(\theta\) and \(\theta'\) the associated classes in \(\operatorname{Ext}^1(\mathbf{M}, \mathbf{N})\) . For \(\theta = \theta'\) , it is necessary and sufficient that there exist an A-homomorphism \(h: \mathbf{R} \to \mathbf{R}'\) making the diagram

\[
\begin{array}{c} \mathsf {R} \\ \mathsf {N} \\ \mathsf {f ^ {\prime}} \end{array} \begin{array}{c} \mathsf {f} \\ \mathsf {h} \\ \mathsf {R ^ {\prime}} \end{array} \begin{array}{c} \mathsf {g} \\ \mathsf {M} \\ \mathsf {g ^ {\prime}} \end{array}
\]

commutative. Such a homomorphism is necessarily an isomorphism.

The condition is sufficient by cor. 1 of prop. 4. If \(\theta = \theta'\) , we have a commutative diagram with exact rows:

\[
\begin{array}{c} 0 \longrightarrow \mathrm{N} \longrightarrow \mathrm{R} \longrightarrow \mathrm{M} \longrightarrow 0 \\ \uparrow_ {1 _ {\mathrm{N}}} \Big | _ {h ^ {\prime}} \Big | _ {1 _ {\mathrm{M}}} \\ 0 \longrightarrow \mathrm{N} \longrightarrow \mathrm{R} ^ {\prime \prime} \longrightarrow \mathrm{M} \longrightarrow 0 \\ \uparrow_ {1 _ {\mathrm{N}}} \Big | _ {h ^ {\prime \prime}} \Big | _ {1 _ {\mathrm{M}}} \\ 0 \longrightarrow \mathrm{N} \longrightarrow \mathrm{R} ^ {\prime} \longrightarrow \mathrm{M} \longrightarrow 0. \end{array}
\]

The morphisms \(h'\) and \(h''\) are isomorphisms by X, p.~7, cor. 3, and \(h = h'' \circ h'^{-1}\) answers the question. The last assertion follows from loc. cit.

Remark 2. --- Theorem 1 gives a description of Ext \(_{A}^{n}\) (M, N) as a set of equivalence classes of exact sequences; it is easy to describe the group law obtained on this set by transport of structure. Indeed, let θ (resp. θ′) be the class of an exact sequence 0 → N \(\xrightarrow{f_{n+1}}\) R \(_{n}\) \(\xrightarrow{f_n}\) \ldots{} → R \(_{1}\) → M → 0 (resp. 0 → N \(\xrightarrow{f_{n+1}'}\) R \(_{n}'\) \(\xrightarrow{f_n'}\) \ldots{} → R \(_{1}'\) → M → 0). Let Δ : M → M ⊕ M and ∇ : N ⊕ N → N be the A-linear maps defined by Δ(x) = (x, x) for x ∈ M and ∇(y, z) = y + z for y, z ∈ N. Consider the map

\[
m: \operatorname{Ext} _ {\mathbf {A}} (\mathbf {M}, \mathbf {N}) \oplus \operatorname{Ext} _ {\mathbf {A}} (\mathbf {M}, \mathbf {N}) \rightarrow \operatorname{Ext} _ {\mathbf {A}} (\mathbf {M} \oplus \mathbf {M}, \mathbf {N} \oplus \mathbf {N})
\]

defined in the remark, p.~121. With the notations of loc. cit., we have \(\nabla \circ i_{\mathbf{N}} = 1_{\mathbf{N}}\) and \(q_{\mathbf{M}} \circ \Delta = 1_{\mathbf{M}}\) , and consequently \(\theta + \theta' = \nabla \circ m(\theta, \theta') \circ \Delta\) . In view of loc. cit. and cor. 3, p.~120, this yields an exact sequence of class \(\theta + \theta'\) : if for example \(n \geqslant 2\) , one may take the sequence

\[
0 \rightarrow \mathrm{N} \rightarrow \mathrm{R} _ {n} ^ {\prime \prime} \rightarrow \mathrm{R} _ {n - 1} \oplus \mathrm{R} _ {n - 1} ^ {\prime} \xrightarrow {f _ {n - 1} \oplus f _ {n - 1} ^ {\prime}} \dots \rightarrow \mathrm{R} _ {2} \oplus \mathrm{R} _ {2} ^ {\prime} \rightarrow \mathrm{R} _ {1} ^ {\prime \prime} \rightarrow \mathrm{M} \rightarrow 0
\]

where \(R_{n}^{\prime\prime}\) is the quotient of \(R_{n} \oplus R_{n}^{\prime}\) by the submodule formed by the pairs

\[
\left(f _ {n + 1} (x), - f _ {n + 1} ^ {\prime} (x)\right) \quad \text { for } x \in \mathbf {N},
\]

and where \(R_{1}^{\prime\prime}=R_{1}\times_{M}R_{1}^{\prime}\)
% semantic-fragments: legacy-input-seam
\relax{}
